A pesar de ser uno de los algoritmos para cifrar información más utilizados en la actualidad, AES no puede usarse aisladamente, pues surgirían varios problemas:

\begin{enumerate}[label={--}]
    \item Si el mensaje a cifrar no tiene un tamaño que sea múltiplo de 128 bits, el último bloque no podrá cifrarse. Para solventar este problema, se utilizan algoritmos de padding, cuya función es rellenar el resto de bytes del bloque. Uno de los que se utilizan es PKCS\#7 \cite[Sección~6.3]{rfc5652}. Consiste en asignar a cada byte restante el número de bytes que faltan para rellenar el bloque, es decir, si el mensaje tiene 5 bytes, faltan 11 bytes para completar el bloque, por lo que se asigna a cada byte restante el elemento 11 en hexadecimal: \texttt{0b}. 
    \item Al ser AES un algoritmo de cifrado por bloques, bloques de texto claro iguales se transforman en bloques cifrados iguales, y en consecuencia se puede obtener información acerca del mensaje original. Para evitar esto, los cifrados por bloques utilizan modos de operación, los cuales definen qué información recibe cada bloque. En la publicación \cite{nist80038a} del NIST se documentan varios modos de operación. Por ejemplo, el modo ECB (Electronic CodeBook Mode) es el modo que cifra cada bloque independientemente (no asegura confidencialidad para bloques iguales), y por lo tanto desaconseja su uso. En contraparte, los modos CBC (Cipher Block Chaining Mode) y CTR (CounTeR Mode) proporcionan aleatoriedad, de forma que bloques iguales producirán resultados distintos.
    \item La elección de la contraseña puede no ser segura si esta no es elegida de manera aleatoria. Como un ordenador es incapaz de producir esta aleatoriedad, se requieren sistemas físicos, cuyos resultados luego son operados para obtener la clave. En \cite{nist800133r2}, NIST especifica cómo debe manejarse la generación de claves. 
\end{enumerate}
\begin{Revision}
    Solventando estas limitaciones, AES es uno de los algoritmos de encriptación más seguros del mundo. Tal es así, que el CNSS (Comité de Sistemas de Seguridad Nacional) de Estados Unidos, lo especifica como apto para cifrar información hasta Top Secret, el estándar más estricto de seguridad, usando en dicho caso claves de 256 bits \cite{cnssp15_2016}. Así pues, a día de hoy no hay ataques conocidos que puedan romper el algoritmo. Los intentos más avanzados consiguen reducir el número de operaciones de AES-128, de $2^{128}$ operaciones que requeriría por fuerza bruta, a $2^{126.1}$; de AES-192, se reduce de $2^{192}$ a $2^{189.7}$ operaciones; y de AES-256 $2^{256}$ a $2^{254.4}$ operaciones \cite{bogdanov2011biclique}. En vista de que parece imposible de romper con la tecnología actual, la mayoría de ataques no están dirigidos al algoritmo en sí, sino a su implementación, por ejemplo, los ataques de canal lateral miden el tiempo que tarda el procesador en realizar las operaciones, las emisiones electromagnéticas que emite o su consumo de energía, para poder así inferir la clave $k$ usada. Es por ello que el foco de la seguridad de AES se pone en realizar implementaciones seguras, que sean robustas ante este tipo de ataques.
\end{Revision}
